% !PRINTED 631
\weberpage{631}
%<<CONT
Eigenschaft hat, dass jede Wurzel einer irreduciblen metacyklischen Gleichung $n^{\mathrm{ten}}$ Grades in $\Omega$ in diesem Ausdruck enthalten ist, und dass auch umgekehrt jeder solche Ausdruck einer irreduciblen Gleichung $n^{\mathrm{ten}}$ Grades in $\Omega$ genügt.

Die Grundlage für diese Untersuchung liefern uns die Resolventen von \gesperrt{Lagrange}.

Es sei jetzt $n$ eine Primzahl, $x_0, x_1, x_2 \ldots x_{n-1}$ sei ein System unabhängiger Variablen, und die Bezeichnung so gewählt, dass $x_n = x_0$, $x_{n+1} = x_1 \ldots$ ist. Endlich sei $\varepsilon$ eine $n^{\mathrm{te}}$ Einheitswurzel. Wir setzen dann
\weq{1}{
(\varepsilon, x) = x_0 + \varepsilon x_1 + \varepsilon^2 x_2 + \cdots + \varepsilon^{n-1} x_{n-1} = \sum_{0, n-1}^{h} \varepsilon^h x_h.
}
Diese Ausdrücke haben wir schon im vierzehnten Abschnitte unter dem Namen der \gesperrt{Resolventen von Lagrange} kennen gelernt, und wir haben dort (\S.\,164) ihr Verhalten gegenüber den \gesperrt{cyklischen Permutationen} untersucht.

Es kommt jetzt darauf an, den Einfluss linearer Permutationen von der im \S.\,180 betrachteten Art auf diese Functionen festzustellen.

Wir schicken einen Hülfssatz voraus, den wir eigentlich nur aus dem Früheren zu reproduciren brauchen. Die \gesperrt{imaginären} $n^{\mathrm{ten}}$ Einheitswurzeln $\varepsilon$ sind die Wurzeln einer Gleichung $X = 0$, wenn
\weq{2}{
X = \xi^{n-1} + \xi^{n-2} + \cdots + \xi + 1
}
gesetzt ist, und $\xi$ eine Variable bedeutet.

Die Function $X$ ist, wie wir früher gesehen haben (\S.\,134), irreducibel im Körper $R$ der rationalen Zahlen.

Legen wir irgend einen Körper $\Omega$ zu Grunde, in dem $X$ irreducibel ist, und leiten daraus den Körper $\Omega(\varepsilon)$ ab, der aus allen rationalen Functionen von $\varepsilon$ mit Coëfficienten in $\Omega$ besteht, so ergiebt sich nach \S.\,142, dass jede Grösse dieses Körpers auf eine und nur auf eine Weise in die Form
\weq{3}{
c_0 + c_1 \varepsilon + c_2 \varepsilon^2 + \cdots + c_{n-2} \varepsilon^{n-2}
}
gesetzt werden kann, worin die $c$ Grössen in $\Omega$ sind.

Der Kürze wegen wollen wir diese Form der Grössen in $\Omega(\varepsilon)$ die \gesperrt{Normalform} nennen und wollen also den Satz aussprechen:

1. \gesperrt{Ist $\Omega$ ein Körper, in dem die Kreistheilungsgleichung $X = 0$ irreducibel ist, so kann jede}
%CONT>>
